\paragraph{The compatible modular representations.}
Choose a number field containing all normalization scalars, the values
of $b$, and a splitting field of $B$. For each coefficient prime $\ell$
and an embedding into $\overline{\mathbb Q}_\ell$, set
\[
 \rho_\ell(\sigma)=b(\sigma)^{-1}a_\sigma\mu_\sigma\circ\sigma.
\]
The cocycle identity makes this a representation commuting with $B$;
after splitting $B$, matrix units give $\rho_\ell=r_\ell\oplus r_\ell$
with $\dim r_\ell=2$. On $G_{L'}$, $\rho_\ell$ is the Tate action of
$A$. Faltings' semisimplicity and endomorphism theorems and $\operatorname{End}^0(A)=B$
show that the multiplicity representation is absolutely irreducible.
Proper-smooth comparison and finite descent, as in
Lemma~\ref{lem:de-rham-descent}, make $r_\ell$ de Rham with Tate
Hodge--Tate weights $\{0,-1\}$; Berger's monodromy theorem makes it
potentially semistable.

It is odd. On Betti homology complex conjugation followed by the
holomorphic isogeny $\mu_c$ exchanges the two Hodge summands, each
a simple $B\otimes\mathbb C$-module. Its action on the multiplicity
space is therefore nonscalar. After the scalar correction its square
is one, so its eigenvalues are $1,-1$. The representation is unramified
outside $T\cup\{\ell\}$, by (LR2) and the reduction and field properties.
Pan's regular two-dimensional modularity theorem applies at $\ell=5$.
We use its published formulation in \cite[Section 7, p.~624]{calegari-icm2022},
with the definition of a geometric representation on pp.~611--612,
and the detailed author proof in \cite[Theorem 1.0.4]{pan2022}.
It requires no additional residual-image restriction at this prime and
identifies $r_5$ with a cuspidal modular representation up to a character
twist. The published formulation is used as stated; no byte or text
identity between Pan's author version and the journal version is asserted.

The twist can be taken finite and the weight is two, without invoking
a separate local-algebraicity theorem for characters. The twisting
character is a de Rham subrepresentation of the relevant Hom
representation. Comparing its integer weight with $\{0,-1\}$ and
the modular weights $\{0,1-k\}$ forces $k=2$ and twist weight zero.
On $G_{L'}$, the identity
$(\det r_5)^2=\det V_5(A)=\chi_{\rm cyc}^2$ shows that
$\det r_5/\chi_{\rm cyc}$ has finite image, since $G_{L'}$ has finite
index in $G_{\mathbb Q}$. The modular determinant
divided by $\chi_{\rm cyc}$ also has finite image. The square of the
twisting character consequently has finite image, so the character
itself has finite image. Absorb it in the modular form and take the
primitive weight-two newform $f$ of level $N$.

The conductor at $5$ is not omitted. For each prime
$p\notin T\cup\{3,5\}$, reduction of $\mu_\sigma$ composed with
relative $p$-Frobenius is a quasi-endomorphism of the good reduction
of $A$ at a chosen place. After clearing denominators its characteristic
polynomial is independent of the Tate prime. Half its trace, multiplied
by $b(\sigma)^{-1}a_\sigma$, is therefore one algebraic number giving
the trace of every $r_\ell$. Enlarge the coefficient field to contain
the coefficients of $f$. The equality with the trace of $f$ at five
is an equality of algebraic numbers; transporting it to an embedding
at three and applying Chebotarev and semisimplicity identifies $r_3$
with the three-adic realization of $f$. Thus use $r_3$ to compute the
conductor at $5$ and $r_5$ at every other prime. The modular level is
the conductor of these realizations by
\cite[\S\S0.3--0.7, Theorem A, pp.~409--411]{carayol}
and newform theory \cite[Theorems 1 and 4]{casselman}.
Here the totally real field is $\mathbb Q$, of odd degree one, so
Carayol's extra finite-place hypothesis for even degree does not apply.
His unramified normalization twist and contragredient preserve the
conductor; his theorem allows arbitrary nebentypus.
Li's Fricke relation and Shimura's
Jacobian construction, as in the baseline modularity argument, also
place $r_5$ in the Tate module of $J_1(N)$.

\paragraph{The primitive level equals the local conductor.}
Let $f\in S_2(\Gamma_0(N),\chi)$ be the primitive newform obtained
after absorbing the finite character twist, and let $\Pi_f$ be its
unitary cuspidal automorphic representation. The character $\chi$
